\begin{textsample}{POS Dim 3 – persona\_gemini – Score 18.00 – t910\_gemini.txt}  \label{ex:f3_pos_043}
The fast \textbf{fashion} \textbf{industry} 's environmental impact is significant, fueled by an obsession with \textbf{consumption}.
Annually, 80 \textbf{billion} garments are produced, with most \textbf{ending} up in landfills.
The resources required are immense ; a single pair of jeans \textbf{uses} 7,000 liters of water, and \textbf{chemical} dyeing contributes to pollution.
This \textbf{growth} is driven by rising \textbf{consumption}, with the \textbf{average} American 's clothing purchases \textbf{increasing} from 34 items in 1991 to 67 in 2007. \textbf{Consumers} are urged to avoid impulse buys and to consider the garment workers behind their \textbf{clothes}, a point highlighted by the anniversary of the Rana Plaza collapse.
While some \textbf{companies} promote donation and \textbf{recycling} efforts, such as H\&M 's \textbf{program}, these are not a \textbf{solution}.
Only 1 % of clothing collected through these \textbf{schemes} is turned into recycled fibers.
Instead of relying on these \textbf{programs}, we can upcycle 95 % of our discarded clothing.
Another \textbf{option} is to simply keep old \textbf{clothes}, as \textbf{trends} often \textbf{cycle} back.
As Vivienne Westwood advises :"Buy less, choose well, make it last."

% matched lemmas: average, billion, chemical, clothes, company, consumer, consumption, cycle, end, fashion, growth, increase, industry, option, program, recycling, scheme, solution, trend, use
\end{textsample}
